\documentclass[header={Lecture 7 Exam Study Notes}]{study-note}

\begin{document}

\StudyTitleBlock{07}{Lecture 7: Exam Study Notes}{Adjoint operators, compact operators, Arzel\`a--Ascoli, and integral operators}

\vspace{0.8em}

\begin{focusbox}
This lecture builds the compact-operator technology used again in Fredholm theory and spectral theory. It also gives the Banach-space adjoint and the key principle that compact operators turn weak convergence into strong convergence.
\end{focusbox}

\tableofcontents

\section{Adjoint Operator}

\begin{defbox}{Adjoint operator in Banach spaces}
If $L\in\mathcal L(X,Y)$, define
\[
L^*:Y^*\to X^*,
\qquad
L^*\psi:=\psi\circ L.
\]
\end{defbox}

\begin{thmbox}{Theorem 9.1: Properties of the adjoint operator}
For $L\in\mathcal L(X,Y)$:
\begin{enumerate}
\item $L^*$ is linear and bounded;
\item $\norm{L^*}=\norm{L}$;
\item $(S\circ L)^*=L^*\circ S^*$;
\item if $L$ is bijective, then $(L^{-1})^*=(L^*)^{-1}$.
\end{enumerate}
\end{thmbox}

\begin{prooflevelbox}{SKETCH}
Know how to verify each property from the definition $L^*\psi=\psi\circ L$. The norm equality $\norm{L^*}=\norm{L}$ uses the dual characterization of the norm.
\end{prooflevelbox}

\section{Compact Operators}

\begin{warningbox}{Common trap}
A compact operator need not be invertible. In an infinite-dimensional Banach space, the identity is not compact.
\end{warningbox}

\begin{defbox}{Compact operator}
An operator $K\in\mathcal L(X,Y)$ is \textbf{compact} if it maps bounded sets in $X$ to relatively compact sets in $Y$.
\end{defbox}

\begin{thmbox}{Theorem 9.2: Norm limits of compact operators are compact}
\thmCompactLimit
\end{thmbox}

\begin{prooflevelbox}{SKETCH}
Approximate the limit operator by compact operators and use a diagonal argument: from any bounded sequence, extract a subsequence whose images under the approximants converge, then pass to the limit.
\end{prooflevelbox}

\begin{thmbox}{Theorem 9.3: Compact operators improve weak convergence}
\thmCompactWeakToStrong
\end{thmbox}

\begin{prooflevelbox}{SKETCH}
Use boundedness of the weakly convergent sequence and total boundedness of the image of the unit ball under a compact operator. The key step: weak convergence gives pointwise convergence on functionals, compactness lets you upgrade to uniform.
\end{prooflevelbox}

\paragraph{Why this matters.}
This is one of the most useful compactness principles in the course: compactness converts weak information into strong information after applying the operator.

\section{Arzel\`a--Ascoli And Kolmogorov--Riesz}

\begin{thmbox}{Arzel\`a--Ascoli theorem}
A subset of $C(K)$ is relatively compact in the sup norm if and only if it is bounded and equicontinuous.
\end{thmbox}

\begin{prooflevelbox}{STATEMENT}
Know the statement and how to apply it. Not the full proof.
\end{prooflevelbox}

\begin{thmbox}{Kolmogorov--Riesz theorem}
For $1\le p<\infty$, a bounded set in $L^p(\R^d)$ is relatively compact if and only if translations are uniformly small and tails are uniformly small.
\end{thmbox}

\begin{prooflevelbox}{STATEMENT}
Know the statement. Not the proof.
\end{prooflevelbox}

\paragraph{Exam use.}
You are expected to know these as compactness criteria, usually without full proofs in oral form.

\section{Schauder And Integral Operators}

\begin{thmbox}{Theorem 9.4: Schauder theorem}
\thmSchauder
\end{thmbox}

\begin{prooflevelbox}{STATEMENT}
$K$ is compact iff $K^*$ is compact. Know the statement; the proof uses double-limit arguments.
\end{prooflevelbox}

\begin{thmbox}{Theorem 9.5: Integral operator with continuous kernel is compact}
\thmIntegralOperator
\end{thmbox}

\begin{prooflevelbox}{STATEMENT}
Know that Arzel\`a--Ascoli is the engine: show boundedness and equicontinuity of the image.
\end{prooflevelbox}

\paragraph{Proof idea.}
Show that $K$ sends the unit ball of $C([a,b])$ to a bounded and equicontinuous family, then apply Arzel\`a--Ascoli.

\section{Exam-Style Exercises}

\begin{exbox}{Exercise 7.1: Finite-rank operators are compact}
Let $X,Y$ be normed spaces and let $K\in\mathcal L(X,Y)$ have finite-dimensional range. Prove that $K$ is compact.

\paragraph{Solution pattern.}
If $(x_n)$ is bounded, then $(Kx_n)$ is bounded in the finite-dimensional space $\operatorname{Ran}K$. Closed bounded sets are compact in finite dimensions, so $(Kx_n)$ has a convergent subsequence.
\end{exbox}

\begin{exbox}{Exercise 7.2: Compactness by Arzel\`a--Ascoli}
Define $T:C([0,1])\to C([0,1])$ by
\[
(Tf)(x)=\int_0^x f(t)\,dt.
\]
Prove that $T$ is compact.

\paragraph{Solution pattern.}
For $\norm{f}_{\infty}\le 1$, show uniform boundedness and equicontinuity:
\[
\abs{(Tf)(x)-(Tf)(y)}\le \abs{x-y}.
\]
Then apply Arzel\`a--Ascoli.
\end{exbox}

\begin{exbox}{Exercise 7.3: Integral operator with continuous kernel}
Let $k\in C([0,1]^2)$ and define
\[
(Kf)(x)=\int_0^1 k(x,y)f(y)\,dy
\]
on $C([0,1])$. Prove that $K$ is bounded and compact.

\paragraph{Solution pattern.}
Boundedness comes from
\[
\norm{Kf}_{\infty}\le \norm{k}_{\infty}\norm{f}_{\infty}.
\]
Compactness again uses Arzel\`a--Ascoli: uniform boundedness is immediate, and equicontinuity follows from uniform continuity of $k$ on the compact square.
\end{exbox}

\begin{exbox}{Exercise 7.4: A noncompact operator}
Show that the identity operator $I:C([0,1])\to C([0,1])$ is not compact.

\paragraph{Solution pattern.}
Use that the closed unit ball of an infinite-dimensional normed space is compact only in finite dimensions, or exhibit a bounded sequence in the unit ball with no uniformly convergent subsequence.
\end{exbox}

\begin{exbox}{Exercise 7.5: Compact operators send weak convergence to strong convergence}
Let $K\in\mathcal L(X,Y)$ be compact and let $x_n\rightharpoonup x$ weakly in $X$. Prove that $Kx_n\to Kx$ strongly in $Y$.

\paragraph{Solution pattern.}
The sequence $(x_n)$ is bounded, so $(Kx_n)$ is relatively compact. Every subsequence has a norm-convergent subsequence; its limit must be $Kx$ because functionals still see the weak limit after composition with $K^*$. Hence the whole sequence converges to $Kx$ in norm.
\end{exbox}

\begin{exbox}{Exercise 7.6: Compute an adjoint}
Let $M_g:L^2(0,1)\to L^2(0,1)$ be multiplication by a bounded function $g\in L^{\infty}(0,1)$. Compute its Hilbert-space adjoint.

\paragraph{Solution pattern.}
Compute the inner product directly:
\[
(M_g f,h)=\int_0^1 g f\overline{h}.
\]
Under the standard convention, this equals
\[
(f,M_{\overline g}h),
\]
so $M_g^*=M_{\overline g}$. For real-valued $g$, the operator is selfadjoint.
\end{exbox}

\begin{exbox}{Exercise 7.7: A finite-sampling operator}
Let $p_0,p_1,p_2\in C([0,1])$ be fixed and define
\[
(Lf)(x)=f(0)p_0(x)+f(1/2)p_1(x)+f(1)p_2(x).
\]
Show that $L:C([0,1])\to C([0,1])$ is bounded and compact.

\paragraph{Solution pattern.}
Estimate by $\norm{f}_{\infty}$ for boundedness. Compactness follows because
\[
\operatorname{Ran}L\subset \operatorname{span}\{p_0,p_1,p_2\},
\]
so the range is finite-dimensional.
\end{exbox}

\section{What To Memorize Precisely}

\begin{focusbox}
Know the exact content of:
\begin{enumerate}
\item Theorems \textbf{9.1}, \textbf{9.2}, \textbf{9.3}, \textbf{9.4}, \textbf{9.5};
\item the definition of compact operator;
\item the statements of Arzel\`a--Ascoli and Kolmogorov--Riesz;
\item the slogan: compact operators map weakly convergent sequences to strongly convergent ones.
\end{enumerate}
\end{focusbox}

\end{document}
